The vectors act on $\Lambda W$ through the Fock representation of \claimref{std-algebra.fock-representation}. (i) For $x = e_S$, $y = e_T$ ($e_S = w_{s_1} \wedge \dots$), $\tau(e_S) \wedge e_T$ is $\pm e_{S \cup T}$ if $S \cap T = \emptyset$ and $0$ otherwise, so $\beta(e_S, e_T) = \pm 1$ exactly when $T$ is the complement of $S$. The Gram matrix is a signed permutation matrix, hence nondegenerate, pairing degree $m$ with $k - m$.

(ii) $\tau$ reverses wedge products, so $\tau(\tau(x) \wedge y) = \tau(y) \wedge x$, and on the top degree $\tau$ acts by $(-1)^{k(k-1)/2}$. So $\beta(y, x) = \operatorname{top}(\tau(\tau(x) \wedge y)) = (-1)^{k(k-1)/2}\beta(x, y)$.

(iii) For $u = w \in W$: $\tau(w \wedge x) = \tau(x) \wedge w$, so $\beta(w x, y) = \operatorname{top}(\tau(x) \wedge w \wedge y) = \beta(x, w y)$. For $u = v \in W'$ acting by the derivation $\iota = \iota_{2B(v, \cdot)}$: on $x \in \Lambda^{m}$, $\iota(\tau x) = (-1)^{m+1} \tau(\iota x)$ (the factor removed from position $j$ of $x$ sits at position $m - j + 1$ of $\tau x$). The top part of $\iota(\tau(x) \wedge y)$ vanishes since $\iota$ lowers degree, and $\iota(\tau x \wedge y) = \iota(\tau x) \wedge y + (-1)^{m} \tau x \wedge \iota y$. Hence $(-1)^{m+1}\beta(\iota x, y) + (-1)^{m}\beta(x, \iota y) = 0$, i.e. $\beta(vx, y) = \beta(x, vy)$. For products, $\beta(u_1 \cdots u_m x, y) = \beta(x, u_m \cdots u_1 y)$, which is $\beta(x, \tilde{a} y)$ for $a = u_1 \cdots u_m$ by \claimref{std-algebra.involution-properties}, and extends linearly. If $\tilde{R} R = 1$ then $\beta(Rx, Ry) = \beta(x, \tilde{R} R y) = \beta(x, y)$.
